% Generated draft: requires mathematical and source review.
\documentclass[11pt]{article}
\usepackage[margin=25mm]{geometry}
\usepackage{fontspec}
\setmainfont{Latin Modern Roman}
\usepackage{amsmath,amssymb,mathtools,unicode-math}
\setmathfont{Latin Modern Math}
\usepackage{enumitem,needspace,xurl,hyperref}
\setlength{\parindent}{0pt}
\setlength{\parskip}{5pt}
\setlength{\emergencystretch}{3em}
\newcommand{\N}{\mathbb{N}}
\newcommand{\Z}{\mathbb{Z}}
\newcommand{\Q}{\mathbb{Q}}
\newcommand{\R}{\mathbb{R}}
\newcommand{\C}{\mathbb{C}}
\newcommand{\F}{\mathbb{F}}
\newcommand{\A}{\mathbb{A}}
\newcommand{\PP}{\mathbb{P}}
\newcommand{\E}{\mathbb{E}}
\newcommand{\eps}{\varepsilon}
\newcommand{\dd}{\,\mathrm{d}}
\begin{document}
\section*{UnsolvedMath ID: 800042}
\textbf{Source problem code:} catalogue entry
\subsection{Definitions and mathematical statement}
A semigroup is a nonempty set equipped with an associative binary operation. An ordinary semigroup variety \(\mathcal V\) is a class of semigroups defined by identities in the language with one binary operation and no distinguished constants. For a nonempty finite set \(X\), the relatively free semigroup \(F_{\mathcal V}(X)\) is a member of \(\mathcal V\), equipped with a generating map \(\iota:X\to F_{\mathcal V}(X)\), such that every map \(f:X\to S\), with \(S\in\mathcal V\), extends uniquely to a semigroup homomorphism \(\bar f:F_{\mathcal V}(X)\to S\) satisfying \(\bar f\circ\iota=f\). Write \(F_{\mathcal V}(n)\) when \(|X|=n\geq 1\).

\textbf{Classification problem attributed to G. Mashevitzky.} Describe the ordinary semigroup varieties \(\mathcal V\) whose finitely generated relatively free semigroups can be separated from each other by identical inclusions.

\textbf{Unresolved formulation.} The supplied record does not define an identical inclusion or its satisfaction relation. It also does not specify whether separation must hold for every pair of distinct positive ranks \(m\neq n\), or only for nonisomorphic relatively free semigroups, nor whether separation has a prescribed direction. These choices affect the quantifiers and the target property, so the classification request cannot be expanded into a fully specified formal predicate from the supplied evidence. The notation above covers positive finite ranks; the source does not specify whether rank zero is intended.

\textbf{Supplied example.} The source states that the variety of semilattices satisfies the requested condition, citing \([\mathrm{M1},\mathrm{M01}]\). Here the semilattice variety consists of semigroups satisfying
\[
xy=yx,\qquad x^2=x.
\]
This example is retained as a source assertion; its intended separation property remains subject to the unresolved terminology above.
\subsection{Short English statement}
The problem asks for a classification of ordinary semigroup varieties whose finitely generated free semigroups can be distinguished by identical inclusions. The supplied source names semilattices as an example, but does not define identical inclusions or specify the exact separation requirement. The registry labels the problem open; that label and the example are supplied information, not independently verified conclusions.
\subsection{Sources}
\begin{itemize}
\item Supplied local UnsolvedMath registry record, permanent ID \(800042\), exact source problem code \texttt{catalogue entry}; title: \emph{The Lyapin's Notebook --- Problem 44}. The record belongs to AMR Open Problem Lists and carries the supplied status \texttt{open}.
\item Bershadsky, Kublanovsky, Mashevitzky, \emph{The Lyapin's Notebook: Unsolved Problems in Semigroup Theory} (2026), Problem 44. Supplied source URL: \url{https://arxiv.org/abs/2604.04763}. The statement attributes the problem specifically to G. Mashevitzky; the registry lists Bershadsky, Kublanovsky, Mashevitzky under proposer metadata and records the year 2026.
\item The supplied provenance records access on 2026-10-16 and extraction by \texttt{arxiv\_latex}. These are supplied dated provenance notes; the source and its mathematical assertions have not been independently checked here. Supplied upstream revision: \texttt{d371ee917ac9647fbb06cd6fc529deec6f3dda82}.
\item References \([\mathrm{M1},\mathrm{M01}]\) are cited in the supplied statement for the semilattice example. Their bibliographic entries and links were not supplied, so they cannot be identified reliably from this record.
\end{itemize}
\end{document}
